\documentclass{scrreprt}

\usepackage[ngerman]{babel}
\usepackage[utf8]{inputenc}
\usepackage[T1]{fontenc}
\usepackage{lmodern}
\usepackage{amsmath, amsthm, amssymb, amscd, amsfonts, amsxtra}
\usepackage{graphicx}

\begin{document}

\begin{titlepage}
\thispagestyle{empty}
\vspace{10mm}
\begin{center}
\includegraphics[height=20mm]{unilogo.eps}
\vfill
\rule{\textwidth}{0.1mm}\vspace{5mm}
\Huge\textbf{Parallele kontinuierliche Kollisionserkennung}
\vspace{1mm}\rule{\textwidth}{0.1mm}
\vfill
\Large
Ausarbeitung\\
Sebastian Brockmeyer und Dominik Krümpelmann
\vfill
\normalsize
Ausarbeitung  im Fachbereich Mathematik/Informatik an der Universität Osnabrück \\
Vorlesung: Parallele Algorithmen bei Henning Wenke

Osnabrück, den \today
\end{center}
\end{titlepage}

\newpage
\mbox{}
\thispagestyle{empty}
\pagenumbering{Roman}
\chapter*{Einleitung}
\addcontentsline{toc}{chapter}{Einleitung}
\setcounter{page}{1}

\newpage
\mbox{}
\thispagestyle{empty}
\newpage
\clearpage
\setcounter{page}{2}
\tableofcontents
\listoffigures
\newpage


%% Beginn des Allgemeinen Teils
\pagenumbering{arabic}

\chapter{Einführung}


\section{Einleitung}
\section{Grundlagen}
\section{Timwarp Algorithmus}
\chapter{Implementations Details}
\section{Paralelliseirung}
\section{Umsetzung}
\chapter{Fazit}
\end{document}